\section{Related Work}
\Secl{related}
Our work extends FVAPPS \cite{dougherty2025proving}, which translated APPS \cite{hendrycks2021apps} coding problems into Lean theorems via inferred PBTs. We apply the same pipeline at scale to PBTs found in the wild, inferring specifications from natural language and source code.

\emph{Lean verification benchmarks.}
Several recent benchmarks target Lean. miniCodeProps \cite{brandfonbrener2024minicodeprops} translates 177 specifications from TIP; CLEVER \cite{thakur2025clever} curates 161 problems from HumanEval; Verina \cite{wu2025verina} covers code, spec, and proof generation across 189 tasks; VeriBench \cite{stechly2025veribench} produces 113 tasks from documented Python; the Vericoding Benchmark \cite{bursuc2025benchmark} spans Lean, Verus, and Dafny with success rates of 27--82\%; and VeriEquivBench \cite{ma2025veriequivbench} uses LeetCode transformations to avoid contamination. All draw specifications from synthetic or pedagogical sources rather than properties written by practicing engineers.

\emph{Other verification tool benchmarks.}
DafnyBench \cite{loughridge2024dafnybench} provides Dafny verification tasks. VerusBench \cite{chen2024verusbench} collects 150 proof tasks spanning Dafny, SV-COMP, and Verus. SV-COMP \cite{beyer2025svcomp} is the standard competition for software verifiers, with over 33{,}000 C tasks, but targets automated verification rather than interactive theorem proving. InvBench \cite{wu2025invbench} focuses specifically on loop invariant synthesis. Chakraborty et al.\ \cite{chakraborty2024neural} explore neural synthesis for SMT-assisted proof-oriented programming in F*.

\emph{Interactive theorem proving and mathematics.}
LeanDojo Benchmark 4 \cite{hsiang2025leandojo} (122{,}517 Mathlib4 theorems) and CoqStoq \cite{thompson2024rango} (196{,}929 theorems from GitHub) are primarily mathematical and largely present in frontier model training data. FrontierMath \cite{glazer2024frontiermath} targets advanced mathematical reasoning; software verification instead demands modelling of program semantics, library APIs, and computational behavior.

\FVSpec differs from prior benchmarks in three ways: (1) specifications originate from real-world PBTs written by practicing engineers; (2) the source PBTs were never formally verified, so most resulting theorems are absent from training data; and (3) we operate at thousand-scale rather than hundreds. This work is motivated by AI safety proposals premised on AI-driven formal verification, including the Guaranteed Safe AI agenda \cite{dalrymple2024guaranteed}.
